\section{How this work was done}\label{app:provenance}

This paper and the work it reports were produced by an AI system, Claude (Anthropic), on 25--27 October 2026. Kaarel H\"anni prompted and steered it by posing the questions of Parts 6--9. No human has reviewed the code, the results or this text.

\paragraph{Organisation.} One orchestrating session took each part's question from the user, wrote protocols, fixed held-out sets, ran evaluations and wrote the study's write-ups. It delegated six kinds of work to separate agents, which reported back to it:
\begin{itemize}
\item design work (\emph{builders});
\item writing the blind held-out toy suite \code{dlheld6} (Part 9);
\item reviews of code and process (\emph{audits});
\item checks of the write-ups against the result files (\emph{claims checks} or \emph{numbers and framing reviews});
\item the first draft of this paper (Parts 6--8);
\item pre-drafting \texttt{PART9.md} and applying the Part 9 claims-check corrections (Part 9).
\end{itemize}
Builders, auditors and the author of \code{dlheld6} had their own instructions and data restrictions; the restrictions on data access are recorded in the protocol files; the builders' were checked by the protocol audits, and the \code{dlheld6} author's from its transcript (its tool calls touched only its own folder and the scratchpad).

\paragraph{Part 6} (25 October; commits on the working branch).
\begin{itemize}
\item \texttt{805bff0e}, 19:59: DFG's harness, a first DFG, and the held-out suite \code{dlheld}.
\item \texttt{359a6b61}, 20:29: DFG v1 frozen for the held-out evaluation.
\item \texttt{0220cf04}, 20:48: v1.1, made post hoc after v1's held-out run, with frozen v1 kept.
\item \texttt{5716239b}, 21:03: the input-permutation wrapper, added after the reproduction audit found the tie-break bias, and before any run on \code{dlheld2}. \code{dlheld2} itself was committed after v1.1 (\texttt{70c604eb}).
\item \texttt{43c842b5}, 21:30: the write-up and all results.
\item \texttt{c12c3b5b}, 21:47: corrections from the numbers and framing reviews.
\end{itemize}
The networks on \code{dlheld} were trained, and their learning curves inspected, before v1 was frozen; DFG was run on that suite only after the freeze. \code{dlheld2} was written by DFG's designer.

\paragraph{Part 7} (25--26 October).
\begin{itemize}
\item \texttt{1ae3b424}, 22:43: datasets, harness, MLP baseline and references.
\item \texttt{ba1f1322}, 23:51: the three frozen adaptations with their dev results, and the post-freeze Cholesky fallback of \code{dfg_gates}.
\item The MLP variants followed, with \code{mlp_l2} added after the fairness audit (\texttt{196a36d1}, \texttt{8c11bf6b}).
\item \texttt{11e7eec7}, 02:01: the write-up.
\item \texttt{9fd619b8}, 02:14: corrections from the claims check.
\end{itemize}
The three builders worked from the dev datasets only, and had about 70 minutes each.

\paragraph{Part 8} (26 October).
\begin{itemize}
\item \texttt{04a215b2}, 06:05: \code{dlheld3} and HELDOUT2 fixed.
\item 06:21--06:27: the shared harness, the protocol and the comparison script were committed. The first builder started at 06:22; the protocol was written before any design result was seen.
\item \texttt{b3489586}, 09:11:20: the three frozen designs with full dev results, and the primary designated. The held-out jobs started at 09:11:27.
\item \texttt{82d7877d} and \texttt{49807a59}, 10:26: the audit-driven additions (permuted toy runs, \code{lin}, network variants with weight decay and logistic loss, \code{mlp_strong}, coverage checks) and the four audit reports.
\item 09:47--16:21: held-out results (the first was \code{hyb_portfolio} on MNIST, \texttt{e5a8cc30}, before the audit commits), references, secondary designs and post-hoc ablations.
\item \texttt{dc51998a}, 16:59: the write-up, after a claims check.
\end{itemize}

\paragraph{What the orchestrator had seen when it changed the evaluation (Part 8).} The four audits ran while the primary's held-out jobs were running. By the time the audit-driven changes were decided, the orchestrating session had seen parts of the primary's held-out results in the job logs:
\begin{itemize}
\item the per-dataset error curves of 7 tabular datasets (phoneme, ring, twonorm and waveform\_40 from the Part 7 set; texture, wine\_quality\_white and sleep from HELDOUT2);
\item the full MNIST learning curve, next to the MLP references then available;
\item a few single-seed cells (connect\_4 $n=30000$, sleep $n=10000$ and 30000);
\item the summaries of the first two \code{random} toy shards, 4 of the 20 targets;
\item the held-out toy results of the \code{hgb} reference.
\end{itemize}
Every change made on the audits' findings strengthens a reference, removes a bias in the learner's favour, or narrows the claim:
\begin{itemize}
\item the network variants;
\item \code{mlp_strong}, with settings from the auditor's recommendation, smoke-tested on dev data only;
\item the input permutation;
\item the restriction of the headline to the fresh sets.
\end{itemize}
No learner was changed after its freeze.

\paragraph{Job scheduling (Part 8).}
\begin{itemize}
\item To give the reference runs priority, the running jobs of \code{hyb_stage} on Fashion-MNIST, led24 and ann\_thyroid were paused with SIGSTOP for about 50 minutes. \code{hyb_stage} has no wall-clock rule, and a stopped process accrues no CPU time, so neither its results nor its recorded CPU time are affected. No \code{hyb_union} job was paused.
\item Two duplicate reference runs, \code{nets_wd} on \code{dlheld} and \code{mlp_strong} on pendigits, were started by mistake and stopped; one copy of each was kept.
\end{itemize}

\paragraph{Part 9} (26--27 October). The protocol, the new development sets and the original held-out sets (\code{dlheld4}, \code{qd_a}, \code{qd_b}) were committed before any Part 9 design work (\texttt{82ca257d}, 20:14 UTC on 26 October), and an amendment fixing the dev image cells followed before design (\texttt{f6ad9afe}, 20:15). The later held-out toy suites came with amendments: \code{dlheld5} (Amendment 2, \texttt{7449313e}, 23:08, after builder S froze) and the Q1 suite \code{dlheld6} (Amendment 3, \texttt{0c4db784}, 02:54 UTC on 27 October, after v2 was frozen).
\begin{itemize}
\item \emph{Builders.} Three independent agents built the sum-feature component (S) and the two image members (I1, I2) on development data and froze before their full dev evaluations. A container restart killed the first I2 session; a new agent resumed from its files (\cref{sec:p9audits}).
\item \emph{Integration and freeze.} The orchestrating session applied the pre-registered rule mechanically from dev result files and froze v2 at 02:11 UTC on 27 October; a one-line crash fix followed before any held-out result existed.
\item \emph{Audits.} Four auditors (protocol, compute, discreteness, fairness) examined code and process (02:13--03:01 UTC on 27 October) without opening any held-out result file (the protocol auditor read the one \code{p9_sum} \code{dlheld5} line the orchestrator had quoted in its chat); the orchestrator's looks at v2's $n\le1000$ progress lines, at 02:40 and 03:03, fell during and just after that window (see below). Their findings led to Amendments 3 and 4 (\cref{tab:p9amend}).
\item \emph{The blind suite.} \code{dlheld6} was written by a fresh agent that was told not to read the design notes, the other suites or any result; its tool calls touched only its own folder and the scratchpad.
\item \emph{What the orchestrator saw.} Before the freeze, three held-out log lines (v1 on \code{qd_a} and \code{qd_b}, and one \code{p9_sum} line on \code{dlheld5}); before Amendment 4, v2's \code{qd_a}/\code{qd_b} test errors at $n\le1000$ in progress logs. Both are recorded in the protocol.
\item \emph{Compute incidents.} Two container restarts; out-of-memory kills of \code{p9_gates}' $n=60000$ dev fits (three times; rerun with the frozen file) and, in the held-out phase, of the resumed baseline \code{mlp} (twice), the \code{mlp_pool} control and one \code{ggbt} ablation, all on \code{qd_a} at $n=60000$, when too many large jobs ran at once (all rerun from saved state, one at a time); and three compute cuts: the CNN context reference was reduced to one network at a fixed weight decay (recorded before it ran on \code{qd_a}/\code{qd_b}) and to two cells on \code{qd_a} (recorded once its $n=1000$ cell had finished, before it ran past that cell), and the CPU-matched MLP ran at $n=60000$ only, at the weight decay \code{mlp_wide} chose (recorded before those runs).
\end{itemize}

\paragraph{This paper.} Its sections on Parts 6--8 were first drafted by a separate agent from the study's write-ups (\texttt{DISCRETE-LEARNER.md}, \texttt{STANDARD-TASKS.md}, \texttt{HYBRID.md} and their result tables), the protocols and the audit reports. The Part 9 section, and the Part 9 revisions of the other sections, were written by the orchestrating session from \texttt{PART9.md}, \texttt{study/part9/RESULTS.md}, the Part 9 protocol and audit reports, and the result files. \texttt{PART9.md}'s methods and process sections were pre-drafted by a separate agent; the corrections from the Part 9 claims check were applied by separate agents, one per group of files, and an independent agent checked each group's edits.
\begin{itemize}
\item The figures (\texttt{make\_figures.py}) and the generated tables (\texttt{make\_tables.py}) are computed directly from the result files with the study's own comparison functions.
\item Their values were compared with the study's markdown tables.
\item An independent claims check of the paper, the README and \texttt{PART9.md} (one checker per group of sections, then an adversarial re-check of each reported discrepancy) compared every claim with the result files, the protocol and the audit reports. Its 69 confirmed findings were applied, the edits were checked by independent verifiers, and their remaining findings were applied by the orchestrating session.
\end{itemize}

\paragraph{Reproducibility.} All code is in \repo{study/}: \texttt{discrete\_learner/} (Part 6), \texttt{standard\_tasks/} (Part 7), \texttt{hybrid/} (Part 8) and \texttt{part9/} (Part 9, with \texttt{PROTOCOL.md}, \texttt{PART9.md}, \texttt{RESULTS.md} and the audit reports). Each write-up ends with the commands that regenerate its results.
\begin{itemize}
\item The networks' raw outputs and the cached datasets are in a session scratch directory, not in the repository. The scripts regenerate them, and \repo{study/hybrid/results/nets_final_err.json} holds the networks' final-step test errors used for the toy comparisons.
\item The orchestration scripts of Part 8 (\texttt{study/hybrid/orchestration/}) are kept on the working branch only.
\item {\sloppy Part 9's job drivers and queue scripts (among them \texttt{run\_v2.sh}, \texttt{jax\_refs.sh}, \texttt{cpu\_refs*.sh}, \texttt{refs5.sh}, \texttt{nets5.sh}, \texttt{after\_restart2.sh}, \texttt{refs\_a4.sh}, \texttt{run\_h6.sh}, \texttt{nets\_r3.sh}, \texttt{final\_q*.sh}), the builders' development scripts and the equivalence checks (\texttt{equiv.py}, \texttt{equiv2.py}) stayed in the session scratch directory and are in neither branch. The drivers only wrapped the committed scripts, and \texttt{PART9.md} gives the commands that reproduce the runs. The Part 6 and 7 audit scripts are likewise in the scratch directory (\cref{app:audits}).\par}
\end{itemize}
